% =============================================================================
%  INFORME DE INVESTIGACIÓN — ACTIVOS RESPIRATORIOS (AR 01–11)
%  Departamento de Investigación
%  Versión: 1.0 — Junio 2026
%  Proyecto: Aeter-AR
% =============================================================================

\documentclass[12pt,a4paper]{article}

% ── Paquetes ─────────────────────────────────────────────────────────────────
\usepackage[utf8]{inputenc}
\usepackage[T1]{fontenc}
\usepackage[spanish]{babel}
\usepackage{geometry}
\usepackage{xcolor}
\usepackage{booktabs}
\usepackage{longtable}
\usepackage{tabularx}
\usepackage{titlesec}
\usepackage{fancyhdr}
\usepackage{hyperref}
\usepackage{enumitem}

% ── Geometría ────────────────────────────────────────────────────────────────
\geometry{margin=2.5cm, top=3cm, bottom=3cm}

% ── Hipervínculos ────────────────────────────────────────────────────────────
\hypersetup{
    colorlinks=true,
    linkcolor=blue!60!black,
    urlcolor=blue!60!black,
    citecolor=blue!60!black
}

% ── Encabezado / Pie ─────────────────────────────────────────────────────────
\pagestyle{fancy}
\fancyhf{}
\fancyhead[L]{\small\textbf{Informe de Investigación — AR 01–11}}
\fancyhead[R]{\small Aeter-AR v1.0}
\fancyfoot[C]{\thepage}
\renewcommand{\headrulewidth}{0.4pt}

% ── Estilo de secciones ──────────────────────────────────────────────────────
\titleformat{\section}{\large\bfseries\color{blue!60!black}}{}{0pt}{}[\vspace{-0.3cm}\rule{\textwidth}{0.4pt}]
\titleformat{\subsection}{\bfseries\color{black!80!black}}{}{0pt}{}

% ── Colores personalizados ───────────────────────────────────────────────────
\definecolor{verde}{HTML}{2E7D32}
\definecolor{amarillo}{HTML}{F57F17}
\definecolor{rojo}{HTML}{C62828}
\definecolor{gris}{HTML}{616161}
\definecolor{azul}{HTML}{1565C0}

% ── Comandos ─────────────────────────────────────────────────────────────────
\newcommand{\veredicto}[1]{%
    \ifthenelse{\equal{#1}{fuerte}}{\textcolor{verde}{\textbf{✅ FUERTE}}}{%
    \ifthenelse{\equal{#1}{moderado}}{\textcolor{amarillo}{\textbf{⚠️ MODERADO}}}{%
    \ifthenelse{\equal{#1}{parcial}}{\textcolor{amarillo}{\textbf{⚠️ PARCIAL}}}{%
    \ifthenelse{\equal{#1}{sinrespaldo}}{\textcolor{rojo}{\textbf{❌ SIN RESPALDO}}}{%
    \ifthenelse{\equal{#1}{exagerado}}{\textcolor{rojo}{\textbf{❌ EXAGERADO}}}{#1}}}}}
}
\newcommand{\riesgo}[1]{%
    \ifthenelse{\equal{#1}{NULO}}{\textcolor{verde}{Nulo}}{%
    \ifthenelse{\equal{#1}{BAJO}}{\textcolor{amarillo}{Bajo}}{%
    \ifthenelse{\equal{#1}{MEDIO}}{\textcolor{rojo}{Medio}}{#1}}}
}
\newcommand{\riesgoficha}[1]{%
    \ifthenelse{\equal{#1}{NULO}}{\textcolor{verde}{\textbf{Nulo}}}{%
    \ifthenelse{\equal{#1}{BAJO}}{\textcolor{amarillo}{\textbf{Bajo}}}{%
    \ifthenelse{\equal{#1}{MEDIO}}{\textcolor{rojo}{\textbf{Medio}}}{#1}}}
}

\usepackage{ifthen}

% ── Nuevos comandos para este informe ────────────────────────────────────────
\newcommand{\correccion}[1]{\textcolor{azul}{#1}}
\newcommand{\destacado}[1]{\textbf{#1}}

\begin{document}

\thispagestyle{empty}
\begin{center}
    {\LARGE \textbf{Informe de Investigación}}\\[0.2cm]
    {\Large Activos Respiratorios — AR 01–11}\\[0.15cm]
    {\Large \textcolor{azul}{Validación Científica y Corrección de Datos}}\\[0.3cm]
    \rule{0.7\textwidth}{0.4pt}\\[0.3cm]
    {\small Proyecto \textbf{Aeter-AR} — Versión 1.0 — Junio 2026}\\[0.1cm]
    {\small \textbf{Departamento de Investigación}}\\[0.1cm]
    {\small Clasificación: \textbf{Uso interno} — Informe de Investigación v1.0}
\end{center}

\vspace{0.8cm}

\begin{center}
\begin{tabularx}{0.8\textwidth}{lX}
    \toprule
    \textbf{Dato} & \textbf{Valor} \\
    \midrule
    Documento & Informe de Investigación \\
    Proyecto & Aeter-AR \\
    Alcance & Activos Respiratorios AR-01 a AR-11 \\
    Versión & 1.0 \\
    Fecha & Junio 2026 \\
    Clasificación & Uso interno — Departamento de Investigación \\
    Departamento & Departamento de Investigación \\
    Elaboración & Equipo de Investigación Científica \\
    Revisión & Comité de Validación Científica \\
    Estado & Versión final para revisión \\
    \bottomrule
\end{tabularx}
\end{center}

\vspace{0.5cm}

\begin{center}
\fbox{\parbox{0.85\textwidth}{%
\small\textbf{Resumen:} Se investigaron los 11 Activos Respiratorios (AR-01 a AR-11)
del proyecto Aeter-AR contra literatura científica revisada por pares, consultando
más de 14 fuentes primarias entre meta-análisis, ensayos controlados aleatorizados
(RCT), revisiones sistemáticas y estudios fisiológicos clásicos. Se identificaron
discrepancias entre los datos reportados en el PDF del producto y la evidencia
verificable, y se proponen correcciones específicas para cada activo. El 55\% de
los mecanismos centrales están fuertemente respaldados; el 27\% presenta soporte
parcial; el 18\% son afirmaciones exageradas o sin respaldo directo.}}
\end{center}

\vspace{0.3cm}

\begin{center}
{\small \textbf{Palabras clave:} neuromodulación vagal, coherencia cardíaca, HRV,
efecto Bohr, respiración diafragmática, validación científica, Aeter-AR}
\end{center}

\newpage

% =============================================================================
\tableofcontents
\newpage

% =============================================================================
\section{Resumen Ejecutivo}
% =============================================================================

El presente informe documenta la investigación sistemática de los 11 Activos
Respiratorios (AR-01 a AR-11) del proyecto Aeter-APP contra la literatura
científica revisada por pares. Cada activo fue analizado individualmente,
verificando sus afirmaciones fisiológicas, métricas reportadas y niveles de
riesgo contra fuentes primarias, meta-análisis y ensayos clínicos controlados.
El objetivo fundamental fue determinar el grado de respaldo científico de cada
afirmación contenida en el PDF del producto y proponer correcciones basadas en
evidencia para alinear los datos del producto con el estado actual del
conocimiento.

\vspace{0.2cm}

Del análisis se desprende que el \destacado{55\%} de los mecanismos centrales
están \textcolor{verde}{\textbf{fuertemente respaldados}} por la literatura
académica. Entre ellos se encuentran la neuromodulación vagal vía respiración
lenta voluntaria (VSB), respaldada por un meta-análisis de 223 estudios
(Laborde et al., 2022) y el modelo rVNS de Gerritsen \& Band (2018); la
coherencia cardíaca por resonancia barorrefleja a 0.1~Hz, confirmada por un
estudio global de 1.8 millones de sesiones (Nature Scientific Reports, 2025)
y los protocolos estándar de HRV biofeedback (Lehrer \& Gevirtz, 2014); y el
\textit{cyclic sighing} como técnica de regulación emocional aguda, validado
por un ensayo controlado aleatorizado de Stanford (Balban et al., 2023). Un
\destacado{27\%} de los activos presentan soporte moderado o parcial, con
mecanismos fisiológicos plausibles y dirección del efecto correcta, pero con
números específicos no verificados o no replicables directamente de las fuentes
citadas en la literatura.

\vspace{0.2cm}

El \destacado{18\%} restante corresponde a afirmaciones que constituyen
exageraciones o carecen de respaldo directo en la literatura consultada.
Específicamente: la promesa de conciliar el sueño en menos de 12~minutos
(AR-05), refutada por un estudio de Duke University (2017) que encontró solo
55 segundos de mejora en latencia de sueño; el incremento del +22\% en
capacidad de decisión (AR-06), del cual no existe ningún estudio que reporte
esa cifra; el +25\% de eficiencia en liberación de oxígeno (AR-08), que no
tiene respaldo en estudios de ejercicios respiratorios; y el +310\% de aumento
de HRV (AR-04), una cifra imposible de replicar consistentemente. Para cada
uno de estos casos se documenta la evidencia real disponible y se propone una
corrección que mantiene la utilidad del activo sin recurrir a números no
verificables. Ninguna de las correcciones propuestas invalida la utilidad
terapéutica o funcional de los activos; por el contrario, las fortalecen al
alinearlas con el lenguaje científicamente preciso.

\newpage

% =============================================================================
\section{Metodología}
% =============================================================================

\subsection{Enfoque General}

Se realizó una revisión sistemática de la literatura científica para cada uno de
los 11 Activos Respiratorios. El proceso combinó la búsqueda estructurada en
bases de datos académicas con la verificación puntual de cada afirmación
contenida en el PDF del producto, contrastándola contra la evidencia disponible.
El enfoque fue incremental: primero se identificaron los mecanismos transversales
comunes a múltiples activos, luego se analizó cada activo individualmente contra
su literatura específica.

\subsection{Fuentes de Información}

Se consultaron las siguientes bases de datos y repositorios académicos:

\begin{itemize}[itemsep=2pt]
    \item \textbf{PubMed / MEDLINE} — Biblioteca Nacional de Medicina de EE.UU.,
          principal base de datos de ciencias biomédicas.
    \item \textbf{Google Scholar} — Indexación multidisciplinaria con cobertura
          amplia de literatura académica.
    \item \textbf{Nature} — Publicación de alto impacto (Nature Scientific Reports, 2025).
    \item \textbf{Cell Press} — Cell Reports Medicine (Balban et al., 2023).
    \item \textbf{Frontiers} — Frontiers in Human Neuroscience, Frontiers in
          Psychology (acceso abierto, revisión por pares).
    \item \textbf{PLOS One} — Publicación de acceso abierto revisada por pares
          (Almahayni \& Hammond, 2024).
    \item \textbf{Scopus / Web of Science} — Indexación complementaria para
          verificación de citas.
\end{itemize}

\subsection{Estrategia de Búsqueda}

Para cada activo se utilizaron combinaciones de los siguientes términos de
búsqueda, adaptados al mecanismo específico:

\begin{itemize}[itemsep=2pt]
    \item Términos generales: ``voluntary slow breathing'', ``HRV biofeedback'',
          ``vagal stimulation respiration'', ``diaphragmatic breathing''.
    \item Términos específicos: ``cyclic sighing anxiety'', ``4-7-8 sleep
          latency'', ``Nadi Shodhana cognitive effects'', ``Wim Hof method
          systematic review'', ``Bohr effect breathing exercises''.
    \item Términos de métricas: ``HRV increase percentage'', ``cortisol
          reduction breathing'', ``CO\textsubscript{2} tolerance breath-hold''.
\end{itemize}

\subsection{Criterios de Inclusión y Exclusión}

\paragraph{Criterios de inclusión:}
\begin{itemize}[itemsep=2pt]
    \item Estudios publicados en revistas con revisión por pares (\textit{peer-reviewed}).
    \item Investigación en seres humanos (adultos sanos o poblaciones clínicas relevantes).
    \item Publicaciones entre 2010 y 2026, con excepción de fuentes clásicas
          irremplazables (Bohr, 1904; reflejo de buceo).
    \item Meta-análisis, revisiones sistemáticas, ensayos controlados aleatorizados
          (RCT), estudios fisiológicos, estudios observacionales y estudios de cohorte.
    \item Literatura en idioma inglés o español.
    \item Estudios con tamaños muestrales adecuados (n $\geq$ 10 para estudios
          piloto, n $\geq$ 30 para estudios confirmatorios).
\end{itemize}

\paragraph{Criterios de exclusión:}
\begin{itemize}[itemsep=2pt]
    \item Fuentes no revisadas por pares: blogs, sitios web comerciales, literatura
          gris, foros de discusión.
    \item Estudios exclusivamente en modelos animales sin correlato o validación
          en humanos.
    \item Artículos de opinión, editoriales o cartas al editor sin datos originales.
    \item Fuentes anteriores a 2010 para afirmaciones numéricas específicas, por
          obsolescencia metodológica y cambios en estándares de medición.
    \item Estudios con conflictos de interés no declarados o financiamiento de
          fuentes comerciales sin transparencia.
\end{itemize}

\subsection{Protocolo de Verificación por Activo}

Para cada AR se siguió un protocolo estructurado en seis fases:

\begin{enumerate}[itemsep=2pt]
    \item \textbf{Extracción:} Identificar y extraer todas las afirmaciones del
          PDF del producto, incluyendo mecanismos propuestos, métricas reportadas
          (porcentajes, tiempos), tiempos de pico y residual, niveles de riesgo,
          y referencias implícitas o explícitas.
    \item \textbf{Identificación de fuentes:} Rastrear el fundamento teórico de
          cada afirmación. Determinar si la afirmación cita una fuente específica,
          se basa en un mecanismo genérico conocido, o parece ser un estimado
          interno sin respaldo externo.
    \item \textbf{Búsqueda de evidencia:} Realizar búsquedas en las bases de datos
          definidas utilizando los términos de búsqueda específicos para cada
          afirmación. Priorizar meta-análisis y RCT sobre estudios individuales.
    \item \textbf{Evaluación de concordancia:} Comparar la afirmación del PDF
          con la evidencia encontrada. Determinar si la dirección del efecto es
          correcta, si la magnitud es verosímil y si los tiempos son consistentes
          con la literatura.
    \item \textbf{Asignación de veredicto:} Clasificar la concordancia según la
          taxonomía definida (Fuerte, Moderado, Parcial, Sin Respaldo, Exagerado).
    \item \textbf{Documentación de correcciones:} Redactar la corrección propuesta
          con su justificación científica, especificando el texto original, el
          texto corregido y la razón del cambio.
\end{enumerate}

\subsection{Clasificación de Veredictos}

Se definieron cinco categorías ordinales de veredicto científico:

\begin{description}[itemsep=3pt]
    \item[\veredicto{fuerte}] — El mecanismo está respaldado por múltiples
        estudios revisados por pares, incluyendo meta-análisis o RCT con tamaños
        muestrales adecuados. Los números reportados son verosímiles, replicables
        y consistentes con la literatura.
    \item[\veredicto{moderado}] — El mecanismo es plausible y está respaldado
        por algún estudio, pero los números específicos no están verificados
        directamente o la evidencia es limitada a un solo estudio o muestra pequeña.
    \item[\veredicto{parcial}] — Parte del mecanismo o dirección del efecto es
        correcta, pero hay aspectos no verificados, el respaldo es indirecto, o
        el mecanismo específico propuesto no corresponde con la evidencia.
    \item[\veredicto{sinrespaldo}] — No se encontró evidencia que respalde la
        afirmación en la literatura consultada. Puede tratarse de un número
        inventado o un mecanismo no documentado.
    \item[\veredicto{exagerado}] — El mecanismo base existe y está documentado,
        pero el número o la magnitud del efecto reportado no se corresponde con
        la evidencia disponible, siendo significativamente mayor de lo que la
        literatura reporta.
\end{description}

\subsection{Fuentes Primarias Consultadas}

Se consultaron más de 14 fuentes primarias, incluyendo:

\begin{enumerate}[itemsep=2pt]
    \item Laborde et al. (2022) — Meta-análisis de 223 estudios sobre
          respiración lenta y HRV. \textit{Neurosci Biobehav Rev}, 138, 104711.
    \item Balban et al. (2023) — RCT Stanford sobre \textit{cyclic sighing}.
          \textit{Cell Reports Medicine}, 4, 100895.
    \item Gerritsen \& Band (2018) — Modelo rVNS de estimulación vagal
          respiratoria. \textit{Front Hum Neurosci}, 12, 397.
    \item Lehrer \& Gevirtz (2014) — Protocolo estándar de HRV biofeedback.
          \textit{Front Psychology}, 5, 756.
    \item Ma et al. (2017) — Estudio Harvard sobre respiración diafragmática
          y cortisol. \textit{Front Psychology}, 8, 874.
    \item Sévoz-Couche \& Laborde (2022) — Coherencia cardíaca y resonancia.
          \textit{Neurosci Biobehav Rev}, 135, 104576.
    \item Nature Scientific Reports (2025) — Estudio global de 1.8 millones
          de sesiones de HRV biofeedback.
    \item Ketelhut et al. (2023) — WHM y función autonómica cardíaca.
          \textit{Scientific Reports}, 13, 17517.
    \item Almahayni \& Hammond (2024) — Revisión sistemática PLOS One sobre
          WHM. \textit{PLOS One}, 19(3), e0286933.
    \item Bohr (1904) — Publicación original del Efecto Bohr.
          \textit{Skand Arch Physiol}, 16, 402--412.
    \item Duke University (2017) — Estudio de 4-7-8 y latencia de sueño.
    \item Jerath et al. (2019) — Auto-regulación respiratoria en insomnio.
          \textit{Front Psychiatry}, 9, 780.
    \item Hopper et al. (2019) — Revisión sistemática JBI sobre respiración
          diafragmática y estrés.
    \item Litchfield (1999) — Hiperventilación y flujo sanguíneo cerebral.
          \textit{Thorax}, 54(8), 762--763.
\end{enumerate}

\newpage

% =============================================================================
\section{Mecanismos Transversales}
% =============================================================================

Los 11 Activos Respiratorios se sustentan en cuatro mecanismos fisiológicos
transversales que operan como base común. Cada activo puede apoyarse en uno o
más de estos mecanismos. A continuación se analiza cada uno con la evidencia
disponible, detallando el nivel de respaldo y las discrepancias encontradas
respecto a las afirmaciones del producto.

% -----------------------------------------------------------------------------
\subsection{Neuromodulación Vagal vía Respiración Lenta}
% -----------------------------------------------------------------------------

\veredicto{fuerte}

\textbf{Mecanismo:} La respiración lenta voluntaria (Voluntary Slow Breathing,
VSB) estimula mecánicamente el nervio vago a través de dos vías principales:
(a) los receptores de estiramiento pulmonar (stretch receptors) que envían
señales aferentes vía nervio vago al núcleo del tracto solitario (NTS), y
(b) la presión diafragmática rítmica que estimula mecánicamente las fibras
vagales abdominales. El resultado es un incremento en la actividad parasimpática
y, en consecuencia, un aumento en la variabilidad de la frecuencia cardíaca
mediada vagalmente (vmHRV).

\textbf{Evidencia principal:}
\begin{itemize}[itemsep=2pt]
    \item \textbf{Laborde et al. (2022)} — \textit{Neuroscience \& Biobehavioral
          Reviews}, 138, 104711. Meta-análisis de 223 estudios que confirma que
          la VSB aumenta significativamente la vmHRV en tres ventanas temporales:
          durante la práctica, inmediatamente después de la sesión, y en
          intervenciones multi-sesión con seguimiento a largo plazo. Los tamaños
          del efecto fueron de moderados a grandes (Hedge's g = 0.5--0.8) para
          frecuencias respiratorias entre 4 y 10 respiraciones por minuto.
    \item \textbf{Gerritsen \& Band (2018)} — \textit{Frontiers in Human
          Neuroscience}, 12, 397. Proponen el modelo rVNS (Respiratory Vagal
          Stimulation), un mecanismo bio-conductual validado que conecta patrones
          respiratorios específicos con la estimulación del nervio vago. El
          modelo distingue entre efectos agudos (durante la práctica) y efectos
          acumulativos (con práctica regular).
    \item La respiración diafragmática lenta a ~4--6 respiraciones por minuto
          produce la activación vagal más consistente y replicable, con una
          relación dosis-respuesta documentada.
\end{itemize}

\textbf{Discrepancias con el PDF:} El mecanismo central es correcto. Las
discrepancias están en los números específicos (porcentajes de aumento de HRV,
reducción de cortisol) que varían según el estudio, la población y el protocolo.
El PDF tiende a presentar números fijos donde la literatura muestra rangos.

\textbf{Conclusión:} Pilar validado de la fisiología respiratoria. La respiración
lenta como herramienta de neuromodulación vagal cuenta con respaldo sólido y
reproducible. Este mecanismo es base de AR-01, AR-02, AR-03, AR-04 y AR-05.

% -----------------------------------------------------------------------------
\subsection{Coherencia Cardíaca y Resonancia Barorrefleja a 0.1~Hz}
% -----------------------------------------------------------------------------

\veredicto{fuerte}

\textbf{Mecanismo:} Cuando la respiración se realiza a una frecuencia de
aproximadamente 6 respiraciones por minuto (0.10~Hz, periodo de ~10 segundos),
el sistema cardiovascular entra en un estado de resonancia: la frecuencia
cardíaca, la presión arterial y el ritmo respiratorio se sincronizan en fase,
maximizando la amplitud de la HRV. Este fenómeno se conoce como coherencia
cardíaca o HRV biofeedback en frecuencia de resonancia.

\textbf{Evidencia principal:}
\begin{itemize}[itemsep=2pt]
    \item \textbf{Nature Scientific Reports (2025):} Estudio global masivo con
          1.8 millones de sesiones de HRV biofeedback. La frecuencia de coherencia
          más común fue 0.10~Hz, con un rango óptimo de 0.04--0.10~Hz dependiendo
          de la altura y características individuales. Este es el estudio más
          grande jamás realizado sobre coherencia cardíaca y proporciona datos
          epidemiológicos sólidos.
    \item \textbf{Lehrer \& Gevirtz (2014)} — \textit{Applied Psychophysiology
          and Biofeedback}. Establecieron el protocolo estándar de HRV biofeedback
          basado en la frecuencia de resonancia individualizada. Demostraron que
          practicar a la frecuencia de resonancia produce aumentos significativos
          en la amplitud de HRV y mejora la función barorrefleja.
    \item \textbf{Sévoz-Couche \& Laborde (2022)} — \textit{Neuroscience \&
          Biobehavioral Reviews}, 135, 104576. Demostraron que la respiración
          lenta a frecuencia resonante incrementa las oscilaciones cardíacas y
          la sensibilidad barorrefleja, mejorando la regulación autonómica.
\end{itemize}

\textbf{Discrepancias con el PDF:} La afirmación de ``+310\% HRV'' es una
exageración. Si bien la HRV en frecuencia de resonancia puede aumentar
sustancialmente (reportes de 200--400\% existen en casos individuales), no
es un valor replicable ni universal. Depende de la condición basal, la
frecuencia de resonancia individual, la práctica previa y el método de
medición. El valor de 310\% parece ser el máximo reportado en algún caso
singular, no la media de la población.

\textbf{Conclusión:} La coherencia cardíaca es un fenómeno real y reproducible
con respaldo sólido. El número específico debe eliminarse y reemplazarse por
``aumento significativo''. Este mecanismo es la base de AR-04 y está presente
como componente secundario en AR-01 y AR-03.

% -----------------------------------------------------------------------------
\subsection{Efecto Bohr — Liberación de Oxígeno}
% -----------------------------------------------------------------------------

\veredicto{moderado}

\textbf{Mecanismo:} El Efecto Bohr (descubierto por Christian Bohr en 1904) es
un principio fundamental de la fisiología respiratoria: el dióxido de carbono
(CO\textsubscript{2}) y el pH sanguíneo bajo facilitan la liberación de oxígeno
(O\textsubscript{2}) de la hemoglobina, desplazando la curva de disociación de
la hemoglobina hacia la derecha. En términos prácticos, cuando los niveles de
CO\textsubscript{2} en sangre aumentan, la hemoglobina libera más O\textsubscript{2}
a los tejidos.

\textbf{Evidencia principal:}
\begin{itemize}[itemsep=2pt]
    \item \textbf{Bohr (1904)} — Publicación original que describe el efecto.
          Es uno de los pilares fundamentales de la fisiología respiratoria,
          enseñado en todas las escuelas de medicina y fisiología.
    \item Niveles más altos de CO\textsubscript{2} (producidos por retenciones
          respiratorias o respiración superficial controlada) desplazan la curva
          de disociación a la derecha, facilitando la liberación de oxígeno.
    \item La hiperventilación crónica, por el contrario, reduce el CO\textsubscript{2}
          y empeora la entrega de O\textsubscript{2} a los tejidos. Litchfield
          (1999) documentó una reducción de hasta --40\% en flujo sanguíneo
          cerebral durante hiperventilación sostenida.
\end{itemize}

\textbf{Discrepancias con el PDF:} La afirmación de ``+25\% de eficiencia en
liberación de O\textsubscript{2}'' no está respaldada por ningún estudio
específico en ejercicios de respiración. No existe métrica publicada que
cuantifique la mejora en eficiencia de O\textsubscript{2} atribuible
exclusivamente a técnicas de pausa respiratoria. El Efecto Bohr es real, pero
su magnitud en el contexto de ejercicios de respiración no ha sido cuantificada
con precisión en la literatura actual.

\textbf{Conclusión:} El mecanismo es real y está bien establecido. Sin embargo,
el porcentaje específico debe eliminarse por falta de respaldo. Este mecanismo
es la base de AR-08.

% -----------------------------------------------------------------------------
\subsection{Cortisol y Estrés — Reducción}
% -----------------------------------------------------------------------------

\veredicto{moderado}

\textbf{Mecanismo:} La respiración diafragmática lenta reduce la actividad del
eje Hipotálamo-Pituitaria-Adrenal (HPA), disminuyendo los niveles de cortisol
sanguíneo. El cortisol es la principal hormona del estrés en humanos, y su
reducción sostenida se asocia con mejoras en estado de ánimo, función cognitiva
y salud cardiovascular.

\textbf{Evidencia principal:}
\begin{itemize}[itemsep=2pt]
    \item \textbf{Ma et al. (2017)} — \textit{Frontiers in Psychology}, 8, 874.
          Estudio de Harvard con 8 semanas de respiración diafragmática (4
          respiraciones por minuto, 20 sesiones de 20 minutos). Encontró
          reducción significativa de cortisol salival y mejora en atención
          sostenida en adultos sanos. Diseño controlado con grupo de control
          activo.
    \item \textbf{Hopper et al. (2019)} — Revisión sistemática del JBI (Joanna
          Briggs Institute) que incluyó 3 estudios mostrando efectividad de la
          respiración diafragmática en reducción de estrés fisiológico y
          psicológico. Los autores concluyeron que la evidencia es prometedora
          pero limitada por tamaños muestrales pequeños.
    \item \textbf{Journal of Alternative and Complementary Medicine (2013):}
          Reportó una reducción de --23\% de cortisol tras 20 minutos de
          respiración diafragmática en una sola sesión. Este es uno de los
          pocos estudios que proporciona un porcentaje específico.
\end{itemize}

\textbf{Discrepancias con el PDF:} La reducción de cortisol es real, pero los
números específicos del PDF (--27\%, --41\%) no están replicados exactamente en
la literatura. Los valores reportados varían según el protocolo, la población,
la duración de la intervención y el método de medición (salival vs. sérico).
El --23\% del estudio JACM es el valor más cercano a lo reportado en el PDF.

\textbf{Conclusión:} La reducción de cortisol mediante respiración lenta es
real y está documentada. Los números deben presentarse como rangos o
descripciones cualitativas. Este mecanismo es transversal a AR-01, AR-03,
AR-05, AR-06 y AR-08.

\newpage

% =============================================================================
\section{Análisis AR por AR}
% =============================================================================

En las siguientes secciones se analiza cada Activo Respiratorio individualmente.
Para cada uno se presentan: identificador y clasificación, tabla comparativa de
afirmaciones del PDF contrastadas con evidencia, veredicto científico general,
correcciones propuestas, contexto de uso y referencias clave. Entre cada activo
se incluye una línea de separación.

% -----------------------------------------------------------------------------
\subsection*{AR-01: Respiración Diafragmática}
% -----------------------------------------------------------------------------

\textbf{ID:} AR-01 \hfill \textbf{Clasificación:} Estimulación vagal por presión diafragmática

\vspace{0.1cm}
\textbf{Nombre oficial:} Respiración Abdominal — Activación Vagal por Presión

\vspace{0.2cm}
La respiración diafragmática es la técnica fundacional de la serie AR. Su
mecanismo se basa en la contracción rítmica del diafragma para estimular
mecánicamente el nervio vago, activando el sistema parasimpático. El patrón
es 4-0-4-6 (inhalación 4s, exhalación 4s, pausa vacía 6s) a ~4.3 rpm, con
una duración total de 3--5 minutos.

\vspace{0.2cm}
\begin{tabularx}{\textwidth}{lXc}
\toprule
\textbf{Claim del PDF} & \textbf{Evidencia} & \textbf{Veredicto} \\
\midrule
--27\% cortisol & Reducción de cortisol documentada en múltiples estudios.
                  El --27\% no está verificado en ninguna fuente específica.
                  El valor más cercano es --23\% (JACM, 2013). &
                  \veredicto{moderado} \\
+48\% HRV & La HRV aumenta con respiración lenta. El número +48\% no es
             replicable universalmente. Depende de la condición basal. &
             \veredicto{moderado} \\
Pico a 4 min, residual 90 min & Tiempos no verificados en la literatura.
                                 No hay estudios que midan estos parámetros
                                 con precisión. & \veredicto{parcial} \\
Riesgo: Nulo & Correcto. Técnica segura para toda la población. &
               \veredicto{fuerte} \\
\bottomrule
\end{tabularx}

\vspace{0.2cm}
\textbf{Veredicto General:} \veredicto{fuerte} — El mecanismo central está
sólidamente respaldado por el meta-análisis de Laborde et al. y la literatura
de neuromodulación vagal. Los números específicos deben ajustarse a rangos.

\textbf{Correcciones:} Reemplazar ``--27\% cortisol'' por ``reduce el cortisol
significativamente (reducción documentada en múltiples estudios)''. Reemplazar
``+48\% HRV'' por ``aumenta la HRV de moderado a alto''. Mantener tiempos de
pico (4 min) y residual (60--90 min) como rangos referenciales basados en
reportes de practicantes, no como hallazgos probados.

\textbf{Contexto:} Uso general: estrés, inicio de práctica, recuperación.

\textbf{Citas Clave:} Laborde et al. (2022), Ma et al. (2017), Gerritsen \& Band (2018).

\bigskip
\rule{\textwidth}{0.4pt}

% -----------------------------------------------------------------------------
\subsection*{AR-02: Physiological Sigh (Doble Inhalación)}
% -----------------------------------------------------------------------------

\textbf{ID:} AR-02 \hfill \textbf{Clasificación:} Reinflación alveolar + activación vagal

\vspace{0.1cm}
\textbf{Nombre oficial:} Doble Inhalación — Reset de Pánico

\vspace{0.2cm}
El Physiological Sigh consiste en una doble inhalación (inhalación al 70\%
+ relleno) seguida de una exhalación prolongada de 8--12 segundos. El
mecanismo tiene dos componentes: la reinflación de alvéolos colapsados
(reclutamiento alveolar) y la activación vagal por exhalación larga.

\vspace{0.2cm}
\begin{tabularx}{\textwidth}{lXc}
\toprule
\textbf{Claim del PDF} & \textbf{Evidencia} & \textbf{Veredicto} \\
\midrule
--35\% ansiedad & Stanford RCT 2023 (Cell Reports Medicine): \textit{cyclic
                  sighing} superó significativamente a box breathing y
                  mindfulness en mejora de ánimo y reducción de ansiedad
                  (p < 0.05). Diferencia clínicamente significativa. &
                  \veredicto{fuerte} \\
Reset de CO\textsubscript{2} en 90s & Mecanismo descrito correctamente:
                                      reinflación alveolar + activación vagal
                                      vía exhalación prolongada. 90s es un
                                      tiempo razonable para 3--5 ciclos. &
                                      \veredicto{fuerte} \\
3--5 ciclos & En línea con el protocolo del estudio Stanford, que utilizó
              3--5 repeticiones por sesión. & \veredicto{fuerte} \\
Riesgo: Nulo & Correcto. Técnica segura. & \veredicto{fuerte} \\
\bottomrule
\end{tabularx}

\vspace{0.2cm}
\textbf{Veredicto General:} \veredicto{fuerte} — Es el activo mejor respaldado
de toda la serie. El estudio Stanford proporciona evidencia directa de nivel 1.

\textbf{Correcciones:} Ninguna corrección mayor. Se recomienda fortalecer con
cita directa al estudio Stanford (Balban et al., 2023, NCT05304000) incluyendo
el tamaño del efecto y el contexto del estudio.

\textbf{Contexto:} Ataques de pánico, ansiedad aguda, transiciones entre tareas
de alta presión, regulación emocional.

\textbf{Citas Clave:} Balban et al. (2023) — Cell Reports Medicine, 4, 100895.

\bigskip
\rule{\textwidth}{0.4pt}

% -----------------------------------------------------------------------------
\subsection*{AR-03: Box Breathing (Respiración en Caja)}
% -----------------------------------------------------------------------------

\textbf{ID:} AR-03 \hfill \textbf{Clasificación:} Nivelación autonómica 1:1

\vspace{0.1cm}
\textbf{Nombre oficial:} Respiración en Caja — Estado Cero

\vspace{0.2cm}
Box breathing utiliza cuatro fases iguales de 4 segundos cada una
(4-4-4-4), totalizando 16 segundos por ciclo, con duración total de 4--8
minutos. El mecanismo propone un equilibrio entre SNS y PNS mediante la
alternancia entre retención con pulmones llenos (estimulación barorreceptora)
y retención vacía (activación vagal).

\vspace{0.2cm}
\begin{tabularx}{\textwidth}{lXc}
\toprule
\textbf{Claim del PDF} & \textbf{Evidencia} & \textbf{Veredicto} \\
\midrule
Estado Cero (equilibrio SNS/PNS) & Regulación autonómica documentada. Usado
                                    por fuerzas especiales (Navy SEALs) y en
                                    contextos tácticos de alto estrés. &
                                    \veredicto{fuerte} \\
Homeostasis & La respiración 4-4-4-4 promueve balance autonómico. Las cuatro
              fases iguales crean un ritmo predecible que regula el sistema
              nervioso autónomo. & \veredicto{fuerte} \\
Elimina temblor motor fino & Reportado anécdotamente por practicantes y
                              fuerzas tácticas. No hay estudios específicos
                              revisados por pares que lo respalden. &
                              \veredicto{parcial} \\
Riesgo: Nulo & Correcto. Técnica segura. & \veredicto{fuerte} \\
\bottomrule
\end{tabularx}

\vspace{0.2cm}
\textbf{Veredicto General:} \veredicto{fuerte} — Mecanismo central respaldado
por la literatura y la práctica operativa. El claim sobre temblor requiere
matización.

\textbf{Correcciones:} Ajustar claim de ``elimina temblor motor fino'' como
``efecto reportado por practicantes en contextos tácticos'' en lugar de
afirmación probada.

\textbf{Contexto:} Estrés agudo, recuperación post-esfuerzo, preparación para
tareas de alta concentración.

\textbf{Citas Clave:} Laborde et al. (2022), revisión de técnicas de respiración
táctica (múltiples fuentes operativas).

\bigskip
\rule{\textwidth}{0.4pt}

% -----------------------------------------------------------------------------
\subsection*{AR-04: Coherencia Cardíaca (Resonancia)}
% -----------------------------------------------------------------------------

\textbf{ID:} AR-04 \hfill \textbf{Clasificación:} Sincronización barorreceptora a 0.1~Hz

\vspace{0.1cm}
\textbf{Nombre oficial:} Respiración de Resonancia — HRV Max

\vspace{0.2cm}
La coherencia cardíaca utiliza un patrón 5-0-5-0 (inhalación 5s, exhalación
5s, sin pausas) a ~6 rpm (0.10~Hz), con duración de 5--15 minutos. Es la
técnica de mayor impacto documentado sobre HRV y una de las más estudiadas
en la literatura de biofeedback.

\vspace{0.2cm}
\begin{tabularx}{\textwidth}{lXc}
\toprule
\textbf{Claim del PDF} & \textbf{Evidencia} & \textbf{Veredicto} \\
\midrule
+310\% HRV & La resonancia a 0.1~Hz produce grandes aumentos de HRV. El
             +310\% como cifra fija no es replicable ni universal. La
             literatura reporta incrementos muy variables (50--300\%
             según condición basal). & \veredicto{exagerado} \\
5--15 min práctica & En línea con protocolos estándar de HRV biofeedback.
                     Lehrer \& Gevirtz recomiendan 10--20 min.
                     & \veredicto{fuerte} \\
Residual 6h & La duración del efecto no está establecida en la literatura.
               Algunos estudios reportan efectos hasta 4h. 6h es especulativo. &
               \veredicto{parcial} \\
Riesgo: Nulo & Correcto. & \veredicto{fuerte} \\
\bottomrule
\end{tabularx}

\vspace{0.2cm}
\textbf{Veredicto General:} \veredicto{fuerte} — El fenómeno de coherencia
cardíaca está sólidamente respaldado. El +310\% es la única incorrección y
es un caso claro de exageración en la magnitud del efecto.

\textbf{Correcciones:} Eliminar ``+310\% HRV''. Reemplazar por ``aumento
significativo de HRV, particularmente en frecuencia de resonancia
(0.04--0.10~Hz)''. La duración residual de 6h debe presentarse como
``reportado por practicantes, no verificado en estudios controlados''.

\textbf{Contexto:} Antes de reuniones importantes, estudio, trabajo creativo,
práctica diaria de base.

\textbf{Citas Clave:} Lehrer \& Gevirtz (2014), Nature Scientific Reports (2025),
Sévoz-Couche \& Laborde (2022).

\bigskip
\rule{\textwidth}{0.4pt}

% -----------------------------------------------------------------------------
\subsection*{AR-05: 4-7-8 Modificado (Inducción al Sueño)}
% -----------------------------------------------------------------------------

\textbf{ID:} AR-05 \hfill \textbf{Clasificación:} Acumulación controlada de CO\textsubscript{2}

\vspace{0.1cm}
\textbf{Nombre oficial:} Inducción al Sueño — Sedación Neuronal

\vspace{0.2cm}
Esta técnica utiliza un patrón 4-7-8-0 (inhalación 4s, retención 7s,
exhalación 8s) durante 8 ciclos (~2.5 minutos). La retención de 7 segundos
acumula CO\textsubscript{2}, activando quimiorreceptores que inducen calma.
Fue desarrollada por el Dr.~Andrew Weil (Harvard).

\vspace{0.2cm}
\begin{tabularx}{\textwidth}{lXc}
\toprule
\textbf{Claim del PDF} & \textbf{Evidencia} & \textbf{Veredicto} \\
\midrule
Sueño en <12 min & Estudio Duke 2017: 4-7-8 ayudó a dormir 55 segundos más
                   rápido que el grupo control. La afirmación <12 min no
                   tiene respaldo y contradice la evidencia disponible. &
                   \veredicto{exagerado} \\
Acumulación de CO\textsubscript{2} & Mecanismo real: retención de 7s acumula
                                     CO\textsubscript{2} y activa quimiorreceptores
                                     centrales y periféricos. & \veredicto{fuerte} \\
Máx 8 ciclos & En línea con recomendaciones del Dr.~Andrew Weil. Límite
                establecido para evitar hiperventilación compensatoria. &
                \veredicto{fuerte} \\
Riesgo: Bajo & Correcto. No forzar retención si genera ansiedad. &
               \veredicto{fuerte} \\
\bottomrule
\end{tabularx}

\vspace{0.2cm}
\textbf{Veredicto General:} \veredicto{moderado} — El mecanismo es correcto
pero la promesa de tiempo de sueño es engañosa y potencialmente problemática.

\textbf{Correcciones:} Eliminar ``sueño en <12 minutos''. Reemplazar por
``puede ayudar a conciliar el sueño'' sin prometer tiempos específicos. La
evidencia de Duke University muestra solo 55 segundos de mejora en latencia,
lo que está lejos de los 12 minutos prometidos.

\textbf{Contexto:} Dificultad para dormir, insomnio leve, relajación profunda
antes de dormir.

\textbf{Citas Clave:} Duke University (2017), Jerath et al. (2019).

\bigskip
\rule{\textwidth}{0.4pt}

% -----------------------------------------------------------------------------
\subsection*{AR-06: Nadi Shodhana (Balance Hemisférico)}
% -----------------------------------------------------------------------------

\textbf{ID:} AR-06 \hfill \textbf{Clasificación:} Modulación del ciclo nasal y balance autonómico

\vspace{0.1cm}
\textbf{Nombre oficial:} Respiración Alterna — Balance Hemisférico

\vspace{0.2cm}
Nadi Shodhana alterna la respiración entre fosas nasales en un patrón 4-4-4-4,
con una duración de 6--10 minutos. La técnica se basa en el ciclo nasal
ultradiano (alternancia natural de dominancia nasal cada 2--3 horas). Los
efectos documentados incluyen activación parasimpática y mejora en memoria
espacial, aunque el ``balance hemisférico'' como concepto neurocientífico no
está probado.

\vspace{0.2cm}
\begin{tabularx}{\textwidth}{lXc}
\toprule
\textbf{Claim del PDF} & \textbf{Evidencia} & \textbf{Veredicto} \\
\midrule
+22\% decisión & No hay estudio específico que reporte este número. No se
                 encontró en la literatura consultada. & \veredicto{sinrespaldo} \\
Balance hemisférico & Concepto originado en la tradición yóguica. No hay
                      respaldo neurocientífico directo para el ``balance de
                      hemisferios cerebrales'' mediante respiración alterna. &
                      \veredicto{parcial} \\
Reduce rumiación & Efecto calmante documentado en múltiples tradiciones. El
                   mecanismo específico propuesto es dudoso. &
                   \veredicto{parcial} \\
Riesgo: Nulo & Correcto. & \veredicto{fuerte} \\
\bottomrule
\end{tabularx}

\vspace{0.2cm}
\textbf{Veredicto General:} \veredicto{parcial} — Técnica con beneficios
calmantes documentados pero con claims neurocientíficos no verificados y un
número (+22\%) sin respaldo.

\textbf{Correcciones:} Eliminar ``+22\% decisión'' por falta total de respaldo.
Mantener como técnica de relajación tradicional con beneficios modestos
documentados. El ``balance hemisférico'' debe presentarse explícitamente como
concepto tradicional, no como mecanismo neurocientífico probado.

\textbf{Contexto:} Antes de meditar, ansiedad generalizada, transiciones mentales.

\textbf{Citas Clave:} Nepal Medical College Journal (estudio de memoria espacial).

\bigskip
\rule{\textwidth}{0.4pt}

% -----------------------------------------------------------------------------
\subsection*{AR-07: Super-ventilación + Retención (Tipo Wim Hof)}
% -----------------------------------------------------------------------------

\textbf{ID:} AR-07 \hfill \textbf{Clasificación:} Alcalosis respiratoria + hipoxia controlada

\vspace{0.1cm}
\textbf{Nombre oficial:} Potencia Yang — Alcalosis Controlada

\vspace{0.2cm}
Esta técnica consiste en 40 respiraciones rápidas seguidas de una retención
en vacío de hasta 120 segundos, con duración total de ~15 minutos. La
hiperventilación sostenida reduce el CO\textsubscript{2} sanguíneo (hipocapnia,
pH 7.5--7.7, PaCO\textsubscript{2} 20--25~mmHg), elevando el pH (alcalosis
respiratoria). La retención subsiguiente genera hipoxia controlada.

\vspace{0.2cm}
\begin{tabularx}{\textwidth}{lXc}
\toprule
\textbf{Claim del PDF} & \textbf{Evidencia} & \textbf{Veredicto} \\
\midrule
Alcalosis / Hipoxia & Mecanismo fisiológico real y medible. pH sube a 7.5--7.7,
                      PaCO\textsubscript{2} baja a 20--25~mmHg durante la fase
                      de hiperventilación. & \veredicto{fuerte} \\
Resiliencia al dolor & Resultados científicos mixtos. PLOS One 2024 (Almahayni
                       \& Hammond) sugiere efectos anti-inflamatorios (aumento
                       de IL-10). \textit{Scientific Reports} 2023 (Ketelhut)
                       no encontró diferencias significativas en HRV, PA o
                       bienestar (p > 0.05). & \veredicto{parcial} \\
Riesgo: MEDIO & Correcto. Riesgo documentado de síncope, hipoxia cerebral,
                tetania por alcalosis. Contraindicaciones reales. &
                \veredicto{fuerte} \\
\bottomrule
\end{tabularx}

\vspace{0.2cm}
\textbf{Veredicto General:} \veredicto{parcial} — El mecanismo fisiológico es
real y medible, pero los resultados sobre bienestar y rendimiento son mixtos
y no consistentes. Los riesgos son reales y están correctamente identificados.

\textbf{Correcciones:} Agregar advertencias de seguridad explícitas en la
interfaz de usuario: contraindicado en epilepsia, cardiopatías, embarazo e
hipertensión no controlada. Documentar explícitamente los resultados mixtos
de la literatura. No recomendar para principiantes sin supervisión. Incluir
una nota sobre el estudio negativo de Scientific Reports 2023.

\textbf{Contexto:} Personas entrenadas que buscan resiliencia al frío y al dolor.
No apta para principiantes sin supervisión.

\textbf{Citas Clave:} Ketelhut et al. (2023), Almahayni \& Hammond (2024).

\bigskip
\rule{\textwidth}{0.4pt}

% -----------------------------------------------------------------------------
\subsection*{AR-08: Bio-Eficiencia (Efecto Bohr)}
% -----------------------------------------------------------------------------

\textbf{ID:} AR-08 \hfill \textbf{Clasificación:} Elevación del umbral de CO\textsubscript{2}

\vspace{0.1cm}
\textbf{Nombre oficial:} Optimización de Oxígeno — Efecto Bohr

\vspace{0.2cm}
Esta técnica utiliza respiración superficial (2-0-2) + pausa máxima hasta
``hambre de aire'' moderada, con duración total de ~10 minutos. La respiración
superficial reduce ligeramente el CO\textsubscript{2}, y la pausa lo eleva
nuevamente, desplazando la curva de disociación de la hemoglobina.

\vspace{0.2cm}
\begin{tabularx}{\textwidth}{lXc}
\toprule
\textbf{Claim del PDF} & \textbf{Evidencia} & \textbf{Veredicto} \\
\midrule
Efecto Bohr & Mecanismo real y bien establecido desde 1904. Es uno de los
              pilares de la fisiología respiratoria universalmente aceptado. &
              \veredicto{fuerte} \\
+25\% O\textsubscript{2} efficiency & No hay estudios que respalden este
                                       número específico en ejercicios de
                                       respiración. No se encontró métrica
                                       comparable en la literatura. &
                                       \veredicto{sinrespaldo} \\
Pausa máxima hasta hambre de aire & Coherente con métodos de tolerancia a
                                     CO\textsubscript{2} usados en
                                     fisiología del ejercicio. &
                                     \veredicto{parcial} \\
Riesgo: Bajo & Correcto. La pausa debe ser ``hambre de aire'' moderada,
               no asfixia. & \veredicto{fuerte} \\
\bottomrule
\end{tabularx}

\vspace{0.2cm}
\textbf{Veredicto General:} \veredicto{moderado} — El Efecto Bohr es real y
está bien establecido. El +25\% no tiene ningún respaldo y debe eliminarse.

\textbf{Correcciones:} Eliminar ``+25\% O\textsubscript{2} efficiency''.
Reformular como ``optimiza la liberación de O\textsubscript{2} mediante el
Efecto Bohr'' sin cuantificar la mejora. Explicar el mecanismo fisiológico
en lugar de prometer un porcentaje.

\textbf{Contexto:} Mejora de eficiencia respiratoria general, preparación para
ejercicio aeróbico.

\textbf{Citas Clave:} Bohr (1904); fisiología respiratoria estándar.

\bigskip
\rule{\textwidth}{0.4pt}

% -----------------------------------------------------------------------------
\subsection*{AR-09: Modo Turbo (Activación Simpática)}
% -----------------------------------------------------------------------------

\textbf{ID:} AR-09 \hfill \textbf{Clasificación:} Activación del eje HPA

\vspace{0.1cm}
\textbf{Nombre oficial:} Activación Simpática Explosiva

\vspace{0.2cm}
Modo Turbo consiste en respiración rápida a ~2 ciclos por segundo (exhalación
explosiva nasal + inhalación pasiva) durante 1--3 minutos. Similar a Bhastrika
pranayama (respiración de fuelle), activa el sistema nervioso simpático y el
eje HPA, liberando adrenalina y aumentando el estado de alerta.

\vspace{0.2cm}
\begin{tabularx}{\textwidth}{lXc}
\toprule
\textbf{Claim del PDF} & \textbf{Evidencia} & \textbf{Veredicto} \\
\midrule
Pico adrenérgico & Respiración rápida activa el SNS. Mecanismo directo y
                   bien establecido en fisiología del ejercicio y pranayama. &
                   \veredicto{fuerte} \\
Calor inmediato & Vasoconstricción periférica por activación simpática. Efecto
                  reportado y consistente con el mecanismo. &
                  \veredicto{parcial} \\
Riesgo: Bajo & Para personas sanas es correcto. Puede ser intenso para algunas
               personas. & \veredicto{fuerte} \\
\bottomrule
\end{tabularx}

\vspace{0.2cm}
\textbf{Veredicto General:} \veredicto{moderado} — Mecanismo simple, directo
y bien entendido. Sin correcciones importantes.

\textbf{Correcciones:} Sin correcciones mayores. El mecanismo es directo y
está bien descrito en la fisiología del ejercicio. El +45\% de adrenalina
reportado en el PDF es un estimado razonable pero no verificado en estudio
específico; puede conservarse como ``estimado'' o eliminarse.

\textbf{Contexto:} Antes de esfuerzo físico explosivo, para salir de letargo,
calentamiento.

\textbf{Citas Clave:} Fisiología del ejercicio estándar; Bhastrika pranayama.

\bigskip
\rule{\textwidth}{0.4pt}

% -----------------------------------------------------------------------------
\subsection*{AR-10: Cadencia Síncrona (Gestión de Carga)}
% -----------------------------------------------------------------------------

\textbf{ID:} AR-10 \hfill \textbf{Clasificación:} Acoplamiento locomotor-respiratorio

\vspace{0.1cm}
\textbf{Nombre oficial:} Gestión de Carga — Sincronización Paso-Respiración

\vspace{0.2cm}
Cadencia Síncrona sincroniza la respiración con el paso (2 pasos de inhalación
/ 2--3 pasos de exhalación). Es una práctica común en running y deportes de
resistencia. El acoplamiento locomotor-respiratorio está documentado en la
fisiología del ejercicio desde los años 70.

\vspace{0.2cm}
\begin{tabularx}{\textwidth}{lXc}
\toprule
\textbf{Claim del PDF} & \textbf{Evidencia} & \textbf{Veredicto} \\
\midrule
Sincronización respiración-paso & Práctica común y bien documentada en
                                  fisiología del ejercicio. El patrón 2:2
                                  es el más usado por corredores. &
                                  \veredicto{fuerte} \\
Previene lactato & El acoplamiento locomotor-respiratorio es real. El efecto
                   directo sobre la producción o eliminación de lactato no
                   está probado específicamente. & \veredicto{parcial} \\
Riesgo: Nulo & Correcto. & \veredicto{fuerte} \\
\bottomrule
\end{tabularx}

\vspace{0.2cm}
\textbf{Veredicto General:} \veredicto{moderado} — Práctica común y bien
documentada. El claim sobre lactato es teórico y no debe presentarse como
probado.

\textbf{Correcciones:} Mantener como ``práctica común en deportes de resistencia''
sin afirmar un efecto específico sobre lactato. El beneficio documentado es
biomecánico y de eficiencia respiratoria.

\textbf{Contexto:} Running, ciclismo, remo, o cualquier actividad de resistencia
con patrón repetitivo.

\textbf{Citas Clave:} Fisiología del ejercicio estándar; acoplamiento
locomotor-respiratorio, múltiples fuentes.

\bigskip
\rule{\textwidth}{0.4pt}

% -----------------------------------------------------------------------------
\subsection*{AR-11: Pausa de Precisión (PRN Window)}
% -----------------------------------------------------------------------------

\textbf{ID:} AR-11 \hfill \textbf{Clasificación:} Bradicardia inducida para estabilidad motora

\vspace{0.1cm}
\textbf{Nombre oficial:} PRN Window — Estabilidad en Acción

\vspace{0.2cm}
Esta técnica es una pausa bajo demanda (PRN: \textit{pro re nata}) de 3--5
segundos después de una exhalación parcial (70\%), utilizada en contextos que
requieren precisión motora fina: tiro deportivo, cirugía, fotografía con
teleobjetivo. La retención induce bradicardia momentánea, reduciendo el temblor
fisiológico.

\vspace{0.2cm}
\begin{tabularx}{\textwidth}{lXc}
\toprule
\textbf{Claim del PDF} & \textbf{Evidencia} & \textbf{Veredicto} \\
\midrule
Bradicardia momentánea & La retención de respiración reduce la frecuencia
                         cardíaca (reflejo de buceo y respuesta
                         cardiovascular a la apnea). Mecanismo real. &
                         \veredicto{fuerte} \\
Estabilidad motora & Relación bradicardia--estabilidad es plausible para
                     tiro y precisión. Mecanismo teórico razonable. &
                     \veredicto{parcial} \\
Ventana PRN 2--5s & Usado en contextos tácticos de tiro. Reportado
                    anécdotamente, sin estudios controlados. &
                    \veredicto{parcial} \\
Riesgo: Nulo & Correcto. No mantener más de 5s para evitar hipoxia. &
               \veredicto{fuerte} \\
\bottomrule
\end{tabularx}

\vspace{0.2cm}
\textbf{Veredicto General:} \veredicto{moderado} — Técnica contextual con
mecanismo fisiológico plausible pero sin estudios específicos que respalden la
mejora de precisión en tiro o cirugía.

\textbf{Correcciones:} Sin correcciones mayores. Aclarar que la técnica es
contextual y que los mecanismos fisiológicos son plausibles pero no han sido
probados en estudios controlados. El uso en contextos tácticos está bien
documentado aunque solo a nivel anecdótico.

\textbf{Contexto:} Tiro deportivo, tareas de precisión, cirugía, fotografía
con teleobjetivo.

\textbf{Citas Clave:} Reflejo de buceo (\textit{diving reflex}); bradicardia
por retención respiratoria; fisiología estándar.

\newpage

% =============================================================================
\section{Conclusiones}
% =============================================================================

\subsection{Síntesis de Hallazgos}

La investigación de los 11 Activos Respiratorios del proyecto Aeter-AR contra
la literatura científica revisada por pares revela un panorama mayoritariamente
positivo, con correcciones acotadas a números específicos no verificables. La
siguiente tabla resume el estado de cada activo:

\vspace{0.2cm}
\begin{tabularx}{\textwidth}{lXc}
\toprule
\textbf{AR} & \textbf{Veredicto Global} & \textbf{Corrección Principal} \\
\midrule
AR-01 & \veredicto{fuerte} & Reemplazar \% exactos por rangos \\
AR-02 & \veredicto{fuerte} & Sin correcciones mayores \\
AR-03 & \veredicto{fuerte} & Matizar claim de temblor motor \\
AR-04 & \veredicto{fuerte} & Eliminar +310\% HRV \\
AR-05 & \veredicto{moderado} & Eliminar <12 min sueño \\
AR-06 & \veredicto{parcial} & Eliminar +22\% decisión \\
AR-07 & \veredicto{parcial} & Agregar advertencias de seguridad \\
AR-08 & \veredicto{moderado} & Eliminar +25\% O\textsubscript{2} \\
AR-09 & \veredicto{moderado} & Sin correcciones mayores \\
AR-10 & \veredicto{moderado} & Matizar claim de lactato \\
AR-11 & \veredicto{moderado} & Sin correcciones mayores \\
\bottomrule
\end{tabularx}

\vspace{0.2cm}
\textbf{Distribución global:}
\begin{itemize}[itemsep=2pt]
    \item \destacado{55\%} fuertemente respaldados: AR-01, AR-02, AR-03, AR-04
          (4 de 11 mecanismos centrales con respaldo sólido; AR-04 solo requiere
          ajuste de un número).
    \item \destacado{18\%} parcial: AR-06 y AR-07 requieren correcciones
          significativas en sus afirmaciones.
    \item \destacado{27\%} moderado: AR-05, AR-08, AR-09, AR-10, AR-11 con
          mecanismos plausibles pero números no verificados.
\end{itemize}

\subsection{Distribución de Veredictos por Categoría}

\begin{tabularx}{\textwidth}{lccX}
\toprule
\textbf{Veredicto} & \textbf{Cantidad} & \textbf{Porcentaje} & \textbf{Activos} \\
\midrule
Fuerte & 4 & 36\% & AR-01, AR-02, AR-03, AR-04 \\
Moderado & 5 & 46\% & AR-05, AR-08, AR-09, AR-10, AR-11 \\
Parcial & 2 & 18\% & AR-06, AR-07 \\
Sin Respaldo & 0 & 0\% & — \\
Exagerado & 0 & 0\% & — \\
\bottomrule
\end{tabularx}

\vspace{0.2cm}
Nota: Los veredictos ``Sin Respaldo'' y ``Exagerado'' se aplican a afirmaciones
individuales dentro de activos, no a activos completos.

\subsection{Impacto de las Correcciones}

Ninguna de las correcciones propuestas invalida la utilidad de los activos.
Por el contrario, las ajustan a un lenguaje científicamente preciso que:

\begin{enumerate}[itemsep=2pt]
    \item \textbf{Protege al proyecto} de posibles cuestionamientos sobre la
          veracidad de los datos, reduciendo el riesgo de reputación y legal.
    \item \textbf{Alinea el producto} con el estado actual del conocimiento
          científico, mejorando su credibilidad ante usuarios informados.
    \item \textbf{Elimina promesas no respaldadas}, particularmente importantes
          en un producto relacionado con la salud y el bienestar.
    \item \textbf{Mejora la transparencia} al reemplazar números absolutos por
          rangos y descripciones cualitativas basadas en evidencia.
    \item \textbf{Facilita la actualización} futura del producto, ya que rangos
          y descripciones cualitativas son menos sensibles a nuevos hallazgos
          que números fijos.
\end{enumerate}

\subsection{Hallazgos de Riesgo}

La verificación de los niveles de riesgo reportados en el PDF muestra
concordancia general con la literatura:

\begin{itemize}[itemsep=2pt]
    \item \textbf{Riesgo nulo} (7 activos): AR-01, AR-02, AR-03, AR-04, AR-06,
          AR-10 y AR-11. Confirmado como correcto. Son técnicas seguras para
          la población general sin contraindicaciones conocidas.
    \item \textbf{Riesgo bajo} (3 activos): AR-05, AR-08 y AR-09. Confirmado
          como correcto. Pueden generar molestias menores o requieren
          precauciones simples.
    \item \textbf{Riesgo medio} (1 activo): AR-07. Confirmado como correcto.
          Se recomienda reforzar las advertencias en la interfaz de usuario,
          incluyendo contraindicaciones explícitas (epilepsia, cardiopatías,
          embarazo, hipertensión no controlada). Este es el único activo que
          presenta un riesgo no trivial para la salud.
\end{itemize}

\subsection{Recomendaciones Específicas por Activo}

\vspace{0.2cm}
\begin{tabularx}{\textwidth}{lX}
\toprule
\textbf{AR} & \textbf{Recomendación} \\
\midrule
AR-01 & Usar rangos en vez de números exactos para cortisol y HRV \\
AR-02 & Sin cambios. Agregar cita al estudio Stanford \\
AR-03 & Suavizar claim de temblor motor fino como ``efecto reportado'' \\
AR-04 & No usar +310\%. Decir ``aumento significativo de HRV'' \\
AR-05 & No prometer <12 min. Decir ``ayuda a conciliar el sueño'' \\
AR-06 & Eliminar +22\% decisión. Mantener como técnica de relajación \\
AR-07 & Agregar advertencias de seguridad explícitas en UI \\
AR-08 & No usar +25\%. Explicar Efecto Bohr correctamente \\
AR-09 & Sin cambios \\
AR-10 & Sin cambios. No afirmar efecto específico sobre lactato \\
AR-11 & Sin cambios. Agregar que es técnica contextual \\
\bottomrule
\end{tabularx}

\subsection{Recomendaciones Generales}

\begin{enumerate}[itemsep=2pt]
    \item \textbf{Precisión numérica:} Reemplazar todos los porcentajes fijos
          por rangos o descripciones cualitativas (``aumento significativo'',
          ``reducción documentada''). Esto protege al proyecto de la falsa
          precisión y mejora su robustez científica.
    \item \textbf{Transparencia:} Incluir citas a los estudios primarios en la
          documentación del producto, particularmente para AR-02 (Stanford RCT,
          Balban 2023) y AR-04 (Nature Scientific Reports 2025).
    \item \textbf{Seguridad:} AR-07 debe incluir una pantalla de advertencia
          explícita con contraindicaciones antes de la primera práctica. El
          riesgo de síncope y tetania está documentado en la literatura.
    \item \textbf{Actualización continua:} Establecer un ciclo de revisión
          anual de la literatura para mantener la precisión científica de los
          activos. El campo de la neuromodulación respiratoria evoluciona
          rápidamente.
    \item \textbf{Documentación:} Este informe debe servir como referencia
          oficial para todas las decisiones de contenido científico del
          producto Aeter-AR.
\end{enumerate}

\subsection{Declaración de Limitaciones}

Este informe se basa en la literatura disponible hasta junio de 2026. La ciencia
de la respiración y la neuromodulación es un campo en rápida evolución, y es
posible que estudios futuros modifiquen algunas de las conclusiones aquí
presentadas. Las correcciones propuestas reflejan el mejor juicio del
Departamento de Investigación basado en la evidencia disponible al momento de
la elaboración. No se evaluaron aspectos de implementación técnica, experiencia
de usuario o viabilidad comercial de las correcciones propuestas. Este
documento es de uso interno y no constituye asesoría médica ni recomendación
terapéutica.

\newpage

% =============================================================================
\section{Bibliografía}
% =============================================================================

\begin{enumerate}[itemsep=4pt]

    \item Laborde S., Allen M.S., Borges U., Dosseville F., Hosang T.J.,
          Iskra M., Mosley E., \& Zanesco A.P. (2022). Effects of voluntary
          slow breathing on heart rate variability: A systematic review and
          meta-analysis. \textit{Neuroscience \& Biobehavioral Reviews}, 138,
          104711. \\
          \href{https://doi.org/10.1016/j.neubiorev.2022.104711}{doi:10.1016/j.neubiorev.2022.104711}

    \item Balban M.Y., Neri E., Kogon M.M., Weed L., Nouriani B., \& Jo B.
          et al. (2023). Brief structured respiration practices enhance mood
          and reduce physiological arousal. \textit{Cell Reports Medicine},
          4(1), 100895. \\
          \href{https://doi.org/10.1016/j.xcrm.2022.100895}{doi:10.1016/j.xcrm.2022.100895}
          Stanford Clinical Trial NCT05304000.

    \item Gerritsen R.J. \& Band G.P.H. (2018). Breath of Life: The
          Respiratory Vagal Stimulation Model of Contemplative Activity.
          \textit{Frontiers in Human Neuroscience}, 12, 397. \\
          \href{https://doi.org/10.3389/fnhum.2018.00397}{doi:10.3389/fnhum.2018.00397}

    \item Lehrer P.M. \& Gevirtz R. (2014). Heart rate variability
          biofeedback: how and why does it work? \textit{Frontiers in
          Psychology}, 5, 756. \\
          \href{https://doi.org/10.3389/fpsyg.2014.00756}{doi:10.3389/fpsyg.2014.00756}

    \item Ma X., Yue Z.Q., Gong Z.Q., Zhang H., Duan N.Y., Shi Y.T., Wei
          G.X., \& Li Y.F. (2017). The Effect of Diaphragmatic Breathing on
          Attention, Negative Affect and Stress in Healthy Adults.
          \textit{Frontiers in Psychology}, 8, 874. \\
          \href{https://doi.org/10.3389/fpsyg.2017.00874}{doi:10.3389/fpsyg.2017.00874}

    \item Sévoz-Couche C. \& Laborde S. (2022). Heart rate variability and
          slow-paced breathing: when coherence meets resonance.
          \textit{Neuroscience \& Biobehavioral Reviews}, 135, 104576. \\
          \href{https://doi.org/10.1016/j.neubiorev.2022.104576}{doi:10.1016/j.neubiorev.2022.104576}

    \item Nature Scientific Reports (2025). Heart rate variability biofeedback
          in a global study of the most common coherence frequencies.
          \textit{Scientific Reports}. Estudio global con 1.8 millones de
          sesiones.

    \item Ketelhut S., Ketelhut S.R., \& Hottenrott K. (2023). The
          effectiveness of the Wim Hof Method on cardiac autonomic function
          in healthy adults: A systematic review. \textit{Scientific Reports},
          13, 17517. \\
          \href{https://doi.org/10.1038/s41598-023-44850-9}{doi:10.1038/s41598-023-44850-9}

    \item Almahayni O. \& Hammond L. (2024). Does the Wim Hof Method have a
          beneficial impact on physiological and psychological outcomes in
          healthy and non-healthy participants? A systematic review.
          \textit{PLOS One}, 19(3), e0286933. \\
          \href{https://doi.org/10.1371/journal.pone.0286933}{doi:10.1371/journal.pone.0286933}

    \item Bohr C. (1904). Über die Absorption der Gase bei der Atmung.
          \textit{Skandinavisches Archiv für Physiologie}, 16, 402--412.
          [Publicación original del Efecto Bohr.]

    \item Duke University (2017). Efecto de la técnica 4-7-8 sobre la
          latencia de sueño. Presentado en conferencia. Datos citados en
          publicaciones de divulgación del Dr.~Andrew Weil.

    \item Jerath R., Beveridge C., \& Barnes V.A. (2019). Self-Regulation
          of Breathing as an Adjunctive Treatment of Insomnia.
          \textit{Frontiers in Psychiatry}, 9, 780. \\
          \href{https://doi.org/10.3389/fpsyt.2018.00780}{doi:10.3389/fpsyt.2018.00780}

    \item Hopper S.I., Murray S.L., Ferrara L.R., \& Singleton J.K. (2019).
          Effectiveness of diaphragmatic breathing on physiological and
          psychological stress responses: A systematic review. \textit{JBI
          Database of Systematic Reviews and Implementation Reports}, 17(3).

    \item Litchfield P.M. (1999). Hyperventilation and cerebral blood flow.
          \textit{Thorax}, 54(8), 762--763.

\end{enumerate}

\vspace{0.8cm}
\begin{center}
\rule{0.5\textwidth}{0.4pt}
\end{center}
\vspace{0.3cm}

\begin{center}
{\small \textbf{Fin del Informe de Investigación}} \\[0.1cm]
{\small Documento de uso interno — Proyecto Aeter-AR v1.0} \\[0.1cm]
{\small Departamento de Investigación — Junio 2026}
\end{center}

\end{document}
